\documentclass[10pt,a4paper]{amsart}
\usepackage[margin=19mm]{geometry}
\usepackage{amsmath,amssymb,mathtools}
\usepackage[scr=rsfs,cal=euler]{mathalfa}
\usepackage{xfrac}
\usepackage[unicode,hidelinks,bookmarks=false]{hyperref}
\newcommand{\ie}{i.e.\ }
\newcommand{\cf}{cf.\ }
\newcommand{\resp}{respectively\ }
\newcommand{\Subsection}{\subsection}
\providecommand{\nobreakdash}{\nobreak\hbox{-}}
\setlength{\emergencystretch}{2em}
\multlinegap0pt
\usepackage{bm}
\usepackage{bbm}
\usepackage{dsfont}
\usepackage{xspace}
\usepackage{cancel}
\usepackage{mathtools}
\usepackage{amsthm}
\makeatletter
\@ifundefined{theorem}{\newtheorem{theorem}{Theorem}}{}
\@ifundefined{assum}{\newtheorem{assum}{Assumption}}{}
\@ifundefined{lemma}{\newtheorem{lemma}{Lemma}}{}
\@ifundefined{proposition}{\newtheorem{proposition}{Proposition}}{}
\makeatother
\newcommand{\ep}{\varepsilon}
\title[Safe Stabilization of the Stefan Problem with a High-Order M]{Safe Stabilization of the Stefan Problem with a High-Order Moving Boundary Dynamics by PDE Backstepping}
\author{Shumon Koga, Miroslav Krstic}
\date{Source: arXiv:2510.06571v1 (CC BY 4.0)}
\begin{document}
\maketitle

\noindent\emph{Source and licence.} Safe Stabilization of the Stefan Problem with a High-Order Moving Boundary Dynamics by PDE Backstepping. Version arXiv:2510.06571v1 (CC BY 4.0), distributed under the Creative Commons Attribution 4.0 International licence, as recorded in the arXiv licence metadata for this exact version and shown on the abstract page https://arxiv.org/abs/2510.06571v1. Full rights notice: licenses/arxiv-2510.06571-CC-BY-4.0.txt.\par\smallskip

\noindent\textit{Editorial scope.} Counted unit: the Proposition whose source label is \texttt{None} in the article (source0.tex lines 904--906, source label \texttt{None}), delivered together with the article's own complete original proof of it (source0.tex lines 907--943, one \texttt{proof} environment of 37 lines). No mathematics is written, repaired or re-derived: statement and proof are the article's own text line for line. Required context retained uncounted: eq:stefanIC (lines 348--355); eq:stefanPDE (lines 348--355); eq:stefanODE (lines 362--366); temp-valid (lines 372--375); ass:initial (lines 379--381); ass:CBFinitial (lines 382--384); lem1 (lines 389--400); eq:qct-valid (lines 391--393); eq:sdot-pos (lines 396--398); eq:qct-exp (lines 526--531); ass:gain-safe (lines 546--556); prop:safety (lines 559--569); eq:st-cond (lines 566--568); eq:h1-def (lines 817--819); eq:h2-ODE (lines 821--824); eq:stefan-thrid-ODE (lines 832--834); eq:third-qct-exp (lines 869--873); ass:a0 (lines 875--877); ass:third-gain12 (lines 878--881); ass:third-setpoint (lines 882--892); ass:third-gain3 (lines 893--901); eq:third-gain3-cond (lines 895--900). Every definition, display and label of those lines is retained unchanged. Omitted from this package and not redistributed: 1--347, 356--361, 367--371, 376--378, 385--388, 394--395, 399--525, 532--545, 557--558, 569--816, 820--820, 825--831, 835--868, 874--874, 901--903, 944--1022. Nothing inside a retained excerpt is omitted or altered: every retained body is delivered exactly as the article's lines are written, so no omission receipt of any kind accompanies this package. Citations: the retained excerpts carry no citation command at all, so no bibliography entry of the article is transcribed and no reference is redistributed. Reference adaptation: none was needed and none was made; every reference inside the delivered unit resolves inside it, so no unresolved reference or citation is produced. Printed slips kept literal: no printed word, symbol, hypothesis or number of the counted statement or of its proof was altered. Layout and class: the article's own class is \texttt{cls/ieeeconf} with packages including xcolor, cite, amsmath, amssymb, amsfonts, amsthm, dsfont, mathtools, nicematrix, adjustbox, graphicx, tabularx, color, multirow; the wrapper is the lane's pinned \texttt{amsart} editorial wrapper at 10pt on A4, which restates the theorem-like environments the retained text needs and adds the article's own macro definitions that the retained text needs (\texttt{\textbackslash ep}), with extra wrapper packages (bm, bbm, dsfont, xspace, cancel, mathtools). The delivered numbering is the unit's own. Assets: no retained span contains a figure environment, a graphics-inclusion command, a PDF-page inclusion, a table or a TikZ picture; no image file is copied into this package, the package declares no asset dependency and no image file is redistributed with it. Licence: version arXiv:2510.06571v1 of this article is distributed under the Creative Commons Attribution 4.0 International licence, as recorded in the arXiv licence metadata for this exact version and shown on the abstract page https://arxiv.org/abs/2510.06571v1. A full rights notice is delivered at licenses/arxiv-2510.06571-CC-BY-4.0.txt. Characters: every retained excerpt is pure ASCII, so the delivered engine, XeLaTeX with its default Latin Modern text font, reports no missing character.
\begin{align}\label{eq:stefanPDE}
T_t(x,t)&=\alpha T_{xx}(x,t), \hspace{2mm}  \textrm{for} \hspace{2mm} t > 0, \hspace{2mm} 0< x< s(t), \\ 
\label{eq:stefancontrol}
-k T_x(0,t)&=q_{\rm c}(t),  \hspace{2mm} \textrm{for} \hspace{2mm} t >0,\\ \label{eq:stefanBC}
T(s(t),t)&= T_{\rm m}, \hspace{2mm} \textrm{for} \hspace{2mm} t >0, \\
\label{eq:stefanIC}
 s(0) &=  s_0, \textrm{and } T(x,0) = T_0(x),  \textrm{for } x \in (0, s_0]. 
\end{align}

\begin{align}
\label{eq:stefanODE}
\ep \ddot s(t) & = - \dot s(t) - \beta T_x(s(t),t) , \\
\dot s(0) & = v_0, 
\end{align}

\begin{align}\label{temp-valid}
T(x,t) \geq& T_{{\rm m}}, \quad  \forall x\in(0,s(t)), \quad \forall t>0, \\
\label{int-valid}0 < s(t)<  &L, \quad \forall t>0. 
\end{align}

\begin{assum}\label{ass:initial} 
$0 < s_0 < L$, $T_0(x) \in C^1([0, s_0];[T_{\rm m}, +\infty))$ with $T_0(s_0) = T_{\rm m}$.
 \end{assum}

 \begin{assum}\label{ass:CBFinitial}
     $v_0 \geq 0$. 
 \end{assum}

 \begin{lemma}\label{lem1}
With Assumptions \ref{ass:initial}-\ref{ass:CBFinitial}, if $q_{\rm c}(t)$ is a bounded piecewise continuous non-negative heat function, i.e.,
\begin{align} \label{eq:qct-valid}
q_{{\rm c}}(t) \geq 0,  \quad \forall t\geq 0,  
\end{align}
then there exists a unique classical solution for the Stefan problem \eqref{eq:stefanPDE}--\eqref{eq:stefanODE}, which satisfies \eqref{temp-valid}, and 
%for all $t \geq 0$, and
\begin{align} \label{eq:sdot-pos} 
    \dot s(t) \geq 0, \quad \forall t \geq 0. 
\end{align}

\end{lemma}

\begin{align}
    q_{\rm c}(t)
    %& = \frac{k}{\alpha}  \int_0^{s(t)} K^\top B u(y,t) dy - k K^\top X(t), \notag\\
    & = - \frac{ k c_2}{\alpha}  \int_0^{s(t)}  (T(x,t) - T_{\rm m})  dx \notag\\
    & \quad - \frac{k}{\beta} \left (c_1 (s(t) - s_{\rm r}) + c_2\ep  \dot s(t) \right).  \label{eq:qct-exp}
\end{align}

\begin{assum} \label{ass:gain-safe}
    The control gains are chosen to satisfy 
    \begin{align} \label{eq:gain-cond}
        c_1 \leq c_2 < \underbrace{c_1 \left(1+\frac{s_{\rm r} - \underline{s}_{\rm r}(s_0, v_0, T_0)}{ \underline{s}_{\rm r}(s_0, v_0, T_0) - s_0}\right)}_{\in(c_1,+\infty]}\,.
        %+ \overline c_2 (s_0, v_0, T_0, s_{\rm r}), 
    \end{align}
    % where 
    % \begin{align}
    %      \overline c_2 (s_0, v_0, T_0, s_{\rm r}) : = c_1\frac{s_{\rm r} - \underline{s}_{\rm r}(s_0, v_0, T_0)}{ \underline{s}_{\rm r}(s_0, v_0, T_0) - s_0}. \label{eq:gain2-bound}    
    % \end{align}
\end{assum}

\begin{proposition} \label{prop:safety}
    Let Assumptions \ref{ass:initial}--\ref{ass:gain-safe} hold. Then, the closed-loop system of \eqref{eq:stefanPDE}--\eqref{eq:stefanODE} with the control law \eqref{eq:qct-exp} 
    % and the conditions of gain parameters 
    % \begin{align} \label{eq:gain-cond}
    %     c_2 \geq c_1, 
    % \end{align}
    has a unique solution satisfying \eqref{temp-valid}, \eqref{eq:qct-valid} and 
    \begin{align} \label{eq:st-cond}
        s_0 \leq s(t) \leq s_{\rm r}, \quad \forall t \geq 0. 
    \end{align}
\end{proposition}

\begin{align} \label{eq:h1-def}
    h_1(t) = \dot s(t).  
\end{align}

\begin{align} \label{eq:h1-ODE}
    \ep_1 \dot h_1(t) &= - h_1 + h_2, \\
 \label{eq:h2-ODE}  \ep_2 \dot h_2(t) &= - h_{2} - \beta T_x(s(t),t), 
\end{align}

\begin{align}
   \ep_1 \ep_2 \dddot s(t) + ( \ep_1 + \ep_2) \ddot s(t) = - \dot s(t) - \beta T_x(s(t),t). \label{eq:stefan-thrid-ODE}
\end{align}

\begin{align}
    q_{\rm c}(t)
    & = - \frac{ k c_3}{\alpha}  \int_0^{s(t)}  (T(x,t) - T_{\rm m})  dx \notag\\
    & \quad - \frac{k}{\beta} \left (c_1 (s(t) - s_{\rm r}) + c_2 (\ep_1 + \ep_2)  \dot s(t) + c_3 \ep_1 \ep_2 \ddot s(t) \right).  \label{eq:third-qct-exp}
\end{align}

\begin{assum} \label{ass:a0}
  It holds that $ a_0 := \ddot s(0) \geq - \frac{v_0}{\ep_1}$. 
\end{assum}

\begin{assum} \label{ass:third-gain12} 
The gains $(c_1, c_2)$ are chosen to satisfy
$0 < c_1 \leq c_2$. 
\end{assum}

\begin{assum} \label{ass:third-setpoint}
The setpoint $s_{\rm r}$ is chosen to satisfy 
    \begin{align}
    s_{\rm r} > \underline{s}_{\rm r}(s_0, v_0, T_0, c_1, c_2)  , \label{eq:third-setpoint}
    \end{align}
    where 
    \begin{eqnarray}
   & \underline{s}_{\rm r}(s_0, v_0, T_0, c_1, c_2) \notag\\
    &:= s_0 + \frac{c_2}{c_1} \left( \ep v_0 + \frac{  \beta}{ \alpha }  \int_0^{s_0}  (T_0(x) - T_{\rm m})  dx \right). \label{eq:third-srunder}
    \end{eqnarray} 
\end{assum}

\begin{assum} \label{ass:third-gain3}
    The gain $c_3$ is chosen to satisfy 
       \begin{align} \label{eq:third-gain3-cond}
        c_2 & \leq c_3 \leq c_2 + \min \left\{\frac{\ep_1}{\ep_2}(c_2 - c_1) , \overline{c}_3, 
        \frac{\ep_2}{\ep_1} c_2
        \right\} , \\
      \overline{c}_3 & :=   \frac{c_1 ( s_{\rm r} - \underline{s}_{\rm r}(s_0, v_0, T_0, c_1, c_2))}{ \ep_1 \ep_2 a_0 +  \frac{ \beta }{\alpha}  \int_0^{s_0}  (T_0(x) - T_{\rm m})  dx }.  
    \end{align}
\end{assum}

\begin{proposition}
    Let Assumptions \ref{ass:initial}, \ref{ass:CBFinitial}, \ref{ass:a0}--\ref{ass:third-gain3} hold. Then, the closed-loop system of \eqref{eq:stefanPDE}--\eqref{eq:stefanIC} and \eqref{eq:stefan-thrid-ODE} under the control law \eqref{eq:third-qct-exp} satisfies \eqref{temp-valid}, \eqref{eq:qct-valid}, \eqref{eq:sdot-pos}, and \eqref{eq:st-cond}. 
\end{proposition}

\begin{proof}
    % Taking the time derivative yields 
% \begin{align}
%     \dot q_{\rm c}(t) &= - c_3 q_{\rm c}(t) \notag\\
%     &+ \frac{k}{\beta} \left( (c_3 - c_1) \dot s(t) + (\ep_1 + \ep_2) (c_3 - c_2) \ddot s(t)\right) \label{eq:third-qcdot}
% \end{align}
Assume $q_{\rm c}(t) > 0 $ for all $t \in [0, t^*)$. Then, due to the CBF \eqref{eq:h2-ODE}, we have $h_2(t) = \ep_1 \ddot s(t) + \dot s(t) \geq 0$. Applying this inequality to the time derivative of \eqref{eq:third-qct-exp}  with the help of the left inequality in \eqref{eq:third-gain3-cond} yields
% that if we have 
% \begin{align} \label{eq:third-gain3}
%  c_3 \geq c_2 , 
% \end{align}
% then \eqref{eq:third-qcdot} yields
\begin{align}
    \dot q_{\rm c}(t) &\geq - c_3 q_{\rm c}(t) \notag\\
    &+ \frac{k}{\beta} \dot s(t) \left( \left(1 + \frac{\ep_2}{\ep_1} \right) c_2 - c_1  - \frac{ \ep_2}{\ep_1} c_3  \right) . \label{eq:third-qctdot-ineq}
\end{align}
Thus, with the first term in the right inequality in \eqref{eq:third-gain3-cond}, 
\eqref{eq:third-qctdot-ineq} yields $\dot q_{\rm c}(t) \geq - c_3 q_{\rm c}(t)$. Note that two satisfy both inequalities of left-hand-side and the first term in the right-hand-side of \eqref{eq:third-gain3-cond}, we need Assumption \ref{ass:third-gain12}. With Assumption \ref{ass:third-setpoint} and the second term in the right-hand side of \eqref{eq:third-gain3-cond}, it holds that $q_{\rm c}(0) > 0$. Applying this and $\dot q_{\rm c}(t) \geq - c_3 q_{\rm c}(t)$ leads to the contradiction of $q_{\rm c}(t^*) = 0$ similarly to the proof of Proposition \ref{prop:safety}. Thus, \eqref{eq:qct-valid} and \eqref{temp-valid} hold by applying Lemma \ref{lem1}, and \eqref{eq:sdot-pos} holds by exponential CBF chain in \eqref{eq:h1-def}--\eqref{eq:h2-ODE} with Assumption \ref{ass:a0}. Finally, applying \eqref{temp-valid}, \eqref{eq:qct-valid}, \eqref{eq:sdot-pos} and $h_2(t) = \ep \ddot s(t) + \dot s(t) >0$ to the control law \eqref{eq:third-qct-exp}, with the third term in the right-hand side of the inequality \eqref{eq:third-gain3-cond}, one can show \eqref{eq:st-cond} holds. 

% leads to the following condition 
% % Another gain condition is imposed as 
% % \begin{align}
% %     \left(1 + \frac{\ep_2}{\ep_1} \right) c_2  > c_1  + \frac{ \ep_2}{\ep_1} c_3 , 
% % \end{align}
% % A comprehensive condition is given as 
% \begin{align} \label{eq:third-gain-cond}
%    c_1 \leq \frac{\ep_1 c_1  + \ep_2 c_3}{\ep_1 + \ep_2 } \leq c_2 \leq c_3, % \quad c_1 \leq c_3 
% \end{align}
% Moreover, we need $q_{\rm c}(0)>0$, and thus substituting $t =0$ in \eqref{eq:third-qct-exp} leads to the condition
% \begin{align}
%    c_3 & <  \frac{c_1 (s_{\rm r} - s_0) - c_2 (\ep_1 + \ep_2)  v_0 }{ \ep_1 \ep_2 a_0 +  \frac{ \beta }{\alpha}  \int_0^{s_0}  (T_0(x) - T_{\rm m})  dx }, \label{eq:third-c3-upper}
% \end{align}
% where $a_0 := \ddot s(0)$. The upper bound of $c_3$ shown in \eqref{eq:third-c3-upper} must be greater than $c_2$, the lower bound of $c_3$ in \eqref{eq:third-gain-cond}, which leads to the setpoint restriction as 
% % Thus the setpoint condition is now given as 

%    Note that all the conditions shown above are for satisfying the safety condition, i.e., $h_1(t) = \dot s(t) \geq 0$ and $h_2(t) = \ep_1 \ddot s(t) + \dot s(t) \geq 0$. 
\end{proof}

\end{document}
